\part{Generadores de números pseudoaleatorios}

\section{Fundamento común}
\label{sec:fundamento}

Los cinco algoritmos que pide el enunciado (LCG, cuadrados medios, Mersenne
Twister, Blum Blum Shub y RANDU) parecen muy distintos entre sí, pero comparten
el mismo esqueleto. Conviene aislarlo una sola vez: de ahí salen dos resultados
—la periodicidad inevitable y la estructura de red— que después se aplican a
cada generador en particular.

\subsection{Qué es, formalmente, un generador pseudoaleatorio}

\begin{definicion}[Generador pseudoaleatorio]
\label{def:prng}
Un \emph{generador de números pseudoaleatorios} es una cuádrupla
$(\mathcal{S}, g, h, x_0)$ donde
\begin{itemize}
  \item $\mathcal{S}$ es un conjunto \textbf{finito} de estados,
  \item $g : \mathcal{S} \to \mathcal{S}$ es la \emph{función de transición},
  \item $h : \mathcal{S} \to [0,1)$ es la \emph{función de salida}, y
  \item $x_0 \in \mathcal{S}$ es la \emph{semilla}.
\end{itemize}
El generador produce la sucesión de estados $x_{n+1} = g(x_n)$ y la sucesión de
salidas $u_n = h(x_n)$.
\end{definicion}

Las tres piezas se reparten así en los algoritmos de esta tarea:

\begin{table}[H]
\centering
\small
\begin{tabularx}{\textwidth}{l L L L}
\toprule
Algoritmo & Estado $\mathcal{S}$ & Transición $g$ & Salida $h$ \\
\midrule
LCG / RANDU & $\{0,\dots,m-1\}$ & $x \mapsto (ax+c) \bmod m$ & $x/m$ \\
Cuadrados medios & enteros de $2k$ dígitos & dígitos centrales de $x^2$ & $x/10^{k}$ \\
Blum Blum Shub & $\mathbb{Z}_M^{*}$ & $x \mapsto x^2 \bmod M$ & bits bajos de $x$ \\
Mersenne Twister & $(\mathbb{Z}_2^{32})^{624}$ & torcedura lineal sobre $\mathbb{F}_2$ & templado de una palabra \\
\bottomrule
\end{tabularx}
\caption{Los cuatro esquemas bajo la misma estructura de la
Definición~\ref{def:prng}.}
\label{tab:esqueleto}
\end{table}

Dos consecuencias se leen directamente de la definición, antes de programar nada:

\paragraph{Determinismo.} Como $g$ es una función y $x_0$ está fijo, toda la
sucesión queda determinada por la semilla. Dos ejecuciones con la misma semilla
producen \emph{exactamente} la misma secuencia. Esto no es un defecto sino el
requisito de reproducibilidad que exige el enunciado; por eso cada generador de
este trabajo recibe su semilla como argumento explícito.

\paragraph{Ausencia de azar genuino.} La sucesión $(u_n)$ no es aleatoria en
ningún sentido probabilístico: es una sucesión fija. Lo único que se puede
pedir es que \emph{sea indistinguible} de una muestra i.i.d.\ $\mathrm{Unif}(0,1)$
frente a una batería de pruebas estadísticas dada. Esta sección se ocupa de la
parte determinista; las pruebas de hipótesis, de la otra.

\subsection{Toda sucesión pseudoaleatoria es periódica}

\begin{teorema}[Periodicidad eventual]
\label{teo:periodicidad}
Sea $(\mathcal{S}, g, h, x_0)$ un generador pseudoaleatorio con
$|\mathcal{S}| = N < \infty$. Entonces existen enteros
$\lambda \ge 0$ (la \emph{cola}) y $\rho \ge 1$ (el \emph{período}) tales que
\[
  x_{n+\rho} = x_n \qquad \text{para todo } n \ge \lambda,
\]
y además $\lambda + \rho \le N$.
\end{teorema}

\begin{proof}
Considérense los $N+1$ estados $x_0, x_1, \dots, x_N$. Todos pertenecen a
$\mathcal{S}$, que tiene solo $N$ elementos. Por el principio del palomar, dos
de ellos coinciden: existen índices $0 \le i < j \le N$ con $x_i = x_j$.

Sea $\rho = j - i \ge 1$. Como $x_{n+1}$ depende \emph{únicamente} de $x_n$,
aplicar $g$ a ambos lados de $x_i = x_j$ da $x_{i+1} = x_{j+1}$, y por inducción
$x_{i+k} = x_{j+k} = x_{i+k+\rho}$ para todo $k \ge 0$. Tomando
$\lambda = i$ se obtiene $x_{n+\rho} = x_n$ para todo $n \ge \lambda$.
Finalmente $\lambda + \rho = i + (j-i) = j \le N$.
\end{proof}

La demostración es corta pero dice tres cosas que importan en todo el trabajo:

\begin{enumerate}
  \item \textbf{El período nunca supera el número de estados}, $\rho \le N$.
        Aumentar el período obliga a agrandar el espacio de estados: es
        exactamente lo que hace el Mersenne Twister al guardar $624$ palabras de
        $32$ bits en lugar de un solo entero, alcanzando $\rho = 2^{19937}-1$.
  \item \textbf{Puede haber una cola $\lambda > 0$}: estados iniciales que nunca
        se vuelven a visitar. Esto ocurre cuando $g$ no es inyectiva, que es
        justamente el caso de los cuadrados medios. Si $g$ \emph{sí} es
        biyectiva —como en un LCG con $\gcd(a,m)=1$— entonces $\lambda = 0$ y la
        sucesión es puramente periódica desde el inicio.
  \item \textbf{La periodicidad es estructural, no un defecto de implementación.}
        Ningún algoritmo determinista sobre memoria finita puede evitarla. Lo
        único negociable es que $\rho$ sea tan grande que no se alcance en la
        práctica.
\end{enumerate}

\subsection{El generador congruencial lineal como caso de estudio}
\label{sec:fundamento-lcg}

El LCG es el más simple de los cinco y el único cuyo comportamiento se puede
caracterizar por completo con álgebra elemental. Por eso se deriva aquí: RANDU
resultará ser un caso particular suyo, y los defectos que se demuestren para la
familia se heredan automáticamente.

La recurrencia es
\begin{equation}
  x_{n+1} = (a\,x_n + c) \bmod m,
  \qquad n = 0, 1, 2, \dots
  \label{eq:lcg}
\end{equation}
con \emph{módulo} $m > 1$, \emph{multiplicador} $0 < a < m$, \emph{incremento}
$0 \le c < m$ y semilla $0 \le x_0 < m$. Si $c = 0$ el generador se llama
\emph{multiplicativo} (o de Lehmer); si $c \neq 0$, \emph{mixto}.

\subsubsection{Forma cerrada}

\begin{lema}[Forma cerrada del LCG]
\label{lema:forma-cerrada}
Para todo $n \ge 0$,
\begin{equation}
  x_n \equiv a^n x_0 + c \sum_{k=0}^{n-1} a^k \pmod{m}.
  \label{eq:forma-cerrada}
\end{equation}
\end{lema}

\begin{proof}
Por inducción sobre $n$.

\emph{Base.} Para $n = 0$ la suma es vacía (vale $0$) y el enunciado se reduce a
$x_0 \equiv a^0 x_0 = x_0$, que es cierto.

\emph{Paso inductivo.} Supóngase \eqref{eq:forma-cerrada} válida para $n$.
Aplicando la recurrencia \eqref{eq:lcg} y la hipótesis de inducción:
\begin{align*}
  x_{n+1} &\equiv a\,x_n + c
   \equiv a\left(a^n x_0 + c \sum_{k=0}^{n-1} a^k\right) + c \\
  &\equiv a^{n+1} x_0 + c \sum_{k=0}^{n-1} a^{k+1} + c
   \equiv a^{n+1} x_0 + c \sum_{k=1}^{n} a^{k} + c \\
  &\equiv a^{n+1} x_0 + c \sum_{k=0}^{n} a^{k} \pmod{m},
\end{align*}
donde el último paso absorbe el término $c = c\,a^0$ dentro de la suma. Esto es
\eqref{eq:forma-cerrada} para $n+1$.
\end{proof}

El Lema~\ref{lema:forma-cerrada} tiene una lectura incómoda: $x_n$ es una
función \emph{explícita y barata} de $n$. No hace falta iterar para saltar a
cualquier punto de la sucesión, y —lo que es peor para cualquier uso
criptográfico— basta observar unas pocas salidas consecutivas para plantear un
sistema y recuperar $(a, c, m)$. Ningún LCG es seguro, y la razón es esta
fórmula, no una debilidad de implementación.

\subsubsection{Período: el teorema de Hull--Dobell}

Por el Teorema~\ref{teo:periodicidad} el período cumple $\rho \le m$. La
pregunta natural es cuándo se alcanza la igualdad.

\begin{teorema}[Hull--Dobell, 1962 \cite{hulldobell}]
\label{teo:hull-dobell}
Un LCG mixto \eqref{eq:lcg} tiene período completo $\rho = m$ para
\textbf{toda} semilla $x_0$ si y solo si se cumplen las tres condiciones:
\begin{enumerate}
  \item $\gcd(c, m) = 1$;
  \item $a - 1$ es divisible por todo factor primo $q$ de $m$;
  \item si $4 \mid m$, entonces $4 \mid (a-1)$.
\end{enumerate}
\end{teorema}

\paragraph{Corolario para módulos potencia de dos.} Si $m = 2^k$ con $k \ge 2$,
el único factor primo de $m$ es $2$ y $4 \mid m$, de modo que las tres
condiciones colapsan en dos reglas verificables de un vistazo:
\begin{equation}
  c \text{ impar}
  \qquad\text{y}\qquad
  a \equiv 1 \pmod 4 .
  \label{eq:hd-potencia2}
\end{equation}
En efecto, $\gcd(c, 2^k) = 1 \iff c$ es impar; la condición (2) pide
$2 \mid (a-1)$ y la (3) pide $4 \mid (a-1)$, de las cuales la segunda implica la
primera. Este corolario explica de entrada por qué RANDU \emph{no} alcanza
período completo: al tener $c = 0$ falla la condición (1) desde el primer paso.

\subsubsection{Los casos degenerados}
\label{sec:degenerados}

Antes de programar conviene identificar qué elecciones de parámetros rompen el
generador, porque son las que la implementación debe rechazar explícitamente.

\paragraph{$a = 0$.} La recurrencia se vuelve $x_{n+1} = c \bmod m$ para todo
$n$: la sucesión es la constante $c/m$ desde el primer término. Período
$\rho = 1$.

\paragraph{$c = 0$ y $x_0 = 0$.} El cero es punto fijo de la recurrencia
multiplicativa, pues $g(0) = (a \cdot 0) \bmod m = 0$. Toda la sucesión es
idénticamente cero. Este caso es especialmente traicionero porque no lanza
ningún error: produce diez mil ceros que pasan desapercibidos si uno no mira el
histograma.

\paragraph{$a = 1$.} La recurrencia se reduce a $x_n = (x_0 + nc) \bmod m$, una
progresión aritmética. Tiene período completo si $\gcd(c,m)=1$, y sin embargo
las salidas son una rampa perfectamente predecible. Es el contraejemplo que
demuestra que \textbf{período completo no implica calidad}: Hull--Dobell es
necesario pero muy lejos de suficiente.

\subsection{Estructura de red: por qué las $k$-tuplas caen en hiperplanos}
\label{sec:red}

El defecto más serio de la familia LCG no lo detecta ninguna prueba
unidimensional. Aparece al mirar las salidas de $k$ en $k$.

\begin{teorema}[Estructura de red \cite{marsaglia}]
\label{teo:red}
Sea un LCG con parámetros $(a, c, m)$ y sea $k \ge 2$. Para todo vector de
enteros $\mathbf{h} = (h_0, h_1, \dots, h_{k-1})$ que satisfaga
\begin{equation}
  h_0 + h_1 a + h_2 a^2 + \dots + h_{k-1} a^{k-1} \equiv 0 \pmod{m},
  \label{eq:condicion-h}
\end{equation}
los puntos $\mathbf{u}_n = (u_n, u_{n+1}, \dots, u_{n+k-1})$ del cubo $[0,1)^k$
caen todos sobre una familia de hiperplanos paralelos, y el número de
hiperplanos efectivamente ocupados es a lo sumo $\sum_{i} |h_i|$.
\end{teorema}

\begin{proof}
Iterando el Lema~\ref{lema:forma-cerrada} desde $x_n$ en lugar de $x_0$,
\[
  x_{n+i} \equiv a^{i} x_n + c\, \sigma_i \pmod m ,
  \qquad
  \sigma_i := \sum_{j=0}^{i-1} a^{j} ,
\]
donde $\sigma_i$ \textbf{no depende de $n$}. Entonces
\[
  \sum_{i=0}^{k-1} h_i\, x_{n+i}
  \;\equiv\; x_n \underbrace{\sum_{i=0}^{k-1} h_i a^{i}}_{\equiv\, 0 \text{ por } \eqref{eq:condicion-h}}
  \;+\; c \underbrace{\sum_{i=0}^{k-1} h_i \sigma_i}_{=:\,K}
  \;\equiv\; cK \pmod m .
\]
El lado izquierdo es congruente módulo $m$ a una constante que no depende de
$n$, así que existe un entero $\ell_n$ con
$\sum_i h_i x_{n+i} = cK + \ell_n m$. Dividiendo entre $m$ y usando
$u_{n+i} = x_{n+i}/m$:
\begin{equation}
  \sum_{i=0}^{k-1} h_i\, u_{n+i} = \frac{cK}{m} + \ell_n .
  \label{eq:hiperplano}
\end{equation}
Para cada valor entero de $\ell_n$, la ecuación \eqref{eq:hiperplano} define un
hiperplano en $[0,1)^k$; al variar $n$ lo único que cambia es $\ell_n$, de modo
que todos los puntos viven en esa familia de hiperplanos paralelos. Como
$u_{n+i} \in [0,1)$, el lado izquierdo de \eqref{eq:hiperplano} está acotado en
un intervalo de longitud $\sum_i |h_i|$, luego $\ell_n$ solo puede tomar a lo
sumo $\sum_i |h_i|$ valores enteros distintos.
\end{proof}

Marsaglia \cite{marsaglia} demuestra además que \emph{siempre} existe un
$\mathbf{h}$ no trivial que cumple \eqref{eq:condicion-h} con coeficientes
pequeños, y que en consecuencia los puntos de cualquier LCG caen en no más de
$(k!\,m)^{1/k}$ hiperplanos. Para $k=3$ y $m = 2^{31}$ esa cota vale
\[
  (3! \cdot 2^{31})^{1/3} = (6 \cdot 2\,147\,483\,648)^{1/3} \approx 2344 ,
\]
que para efectos prácticos es inofensivo: $2344$ planos llenan el cubo de forma
razonablemente densa. El escándalo de RANDU —que se documenta en su propia
sección— es que se queda en \textbf{15}, dos órdenes de magnitud por debajo de
lo que la cota permitía.

\paragraph{Lectura operativa.} El Teorema~\ref{teo:red} explica por qué el
enunciado pide la \emph{prueba espectral visual} y no solo pruebas de hipótesis.
Las pruebas de uniformidad ($\chi^2$, Kolmogórov--Smirnov) miran una coordenada
a la vez; las de independencia serial (autocorrelación, rachas) miran
correlaciones lineales de a pares. Una relación exacta de tres términos como
\eqref{eq:condicion-h} es invisible para todas ellas y salta a la vista en el
diagrama de dispersión tridimensional. Es la diferencia entre no encontrar
evidencia y no haberla buscado donde estaba.

\subsection{Normalización}

Las salidas se obtienen dividiendo el estado entre el módulo, $u_n = x_n / m$.
Se divide entre $m$ y no entre $m-1$ porque el estado recorre
$\{0, 1, \dots, m-1\}$: el valor máximo posible es $(m-1)/m < 1$, que es
exactamente el soporte $[0,1)$ de la distribución $\mathrm{Unif}(0,1)$ estándar.
Dividir entre $m-1$ permitiría la salida $u = 1$, fuera del soporte, lo cual
rompe cualquier transformación posterior que use $\log(1-u)$ o la inversa de una
función de distribución acumulada.

Conviene tener presente que la salida es un valor de una rejilla discreta de
paso $1/m$, no un real del continuo. Ninguna implementación puede producir más
de $m$ valores distintos, y por eso una prueba de bondad de ajuste con más de
$m$ clases carecería de sentido.
